% §2 Preliminaries and notation --- blueprint chapter 2 (blueprint/src/parts/02-preliminaries.tex),
% statements verbatim, proofs of record verbatim. Drafted 2026-09-11, the first section written:
% it fixes the notation table the spatial modules share (hub RELEASES.md rule 3).
%
% Contents and departures from the blueprint chapter:
%   - the letters paragraph and Table 1 are paper-only: the draft's §2 table (the map to the
%     causal article's letters), extended by the symbols this section defines;
%   - all nine statement nodes of the chapter are transcribed, each under its
%     `% shared with blueprint` marker with the blueprint's own label; the [A] nodes' page
%     anchors are printed in the prose before the statement, read from blueprint/AXIOMS.md
%     (A1-A4 at prop:fourier-toolbox, A7 at prop:sd-exponents), never from memory;
%   - the blueprint's status annotations are not transcribed; what a reader needs from them ---
%     which clauses the machine-checked development actually spends, and what a statement does
%     not claim --- is said in prose next to the statement;
%   - rem:levy-continuity-general and rem:displacement-dictionary are copied from the blueprint
%     (remarks carry no marker: they are not statement nodes); "classes referred to later" is
%     the draft's closing paragraph, with its citations;
%   - the causal article appears in Table 1's third column and in exactly one remark,
%     rem:causal-relation-preliminaries, per the reference policy (repo CLAUDE.md);
%   - an equation is cited with \ref everywhere, as the blueprint does (\eqref inside a shared
%     statement would break the byte-identity); macros.tex wraps \ref so that an `eq:` label
%     prints in parentheses, which is what \eqref would have printed;
%   - the blueprint's statements write "\Levy" without the accent, for the render pipeline; the
%     paper keeps them byte-identical (state `verbatim` in `linkage check`). Restoring the
%     accent everywhere is a one-macro change in both macros files, to be decided by the author;
%   - prop:laplace-uniqueness-locally-finite (a chapter-2 node) is NOT transcribed: its consumers
%     are in chapters 10 and 13, outside v0.1, so it left with them on 2026-09-12; "classes
%     referred to later" is trimmed to the two classes v0.1 refers to; lem:cin-rays, which the
%     dictionary remark names, is Appendix A.

\section{Preliminaries and notation}\label{sec:preliminaries}

This section fixes the notation and states the Fourier facts and the profile representation
that the later sections use. The axioms themselves begin in \S\ref{sec:axioms}.

\emph{Letters.} Space runs over $\RR$ with coordinate $x$. The transform variable is the
wavenumber $\omega$. Scale is $t$, the evolution variable of the diffusion equation in the
scale-space literature. At first $t$ is only an ordered label of the records; in the
canonical parametrization of \S\ref{sec:covariance} it has the dimension of length, and the
Gaussian member's diffusion parameter is proportional to $t^2$. A stage is indexed by an ordered pair of scales $s \le t$, with a
third letter $r \le s \le t$ when three are needed. This is the convention for evolution
families, $U(t,s)\,U(s,r) = U(t,r)$, and it matches the indexing of additive processes: the
kernel $\mu_{s,t}$ is the law of an increment $X_t - X_s$. The dilation ratio is $\lambda$.

The causal article \sscite{fagerstrom2026hemigroup} used other letters: the time coordinate
$t$, the Laplace variable $s$, scales $x \le y \le z$ and the ratio $\sigma$.
Table~\ref{tab:notation} records the correspondence, together with the symbols defined in
this section, and it is used whenever a statement of that article is cited. No formula holds
both meanings of $s$ or $t$ at once, since the Laplace variable and the time coordinate do
not occur here.

Two letters carry two roles each, always locally. The letter $a$ is the Gaussian coefficient
from \ref{eq:levy-khintchine} on, and the amount of a translation in $\Tr{a}$. The letter
$c$ is the contraction factor of Definition~\ref{def:self-decomposable}, the orbit coordinate
$c(\lambda)$ of \S\ref{sec:covariance}, and a generic constant.

\begin{table}[!htb]
\centering
\caption{Notation: the symbol, where it is defined, and the letter the causal article uses
for the same object.}
\label{tab:notation}
\footnotesize
\setlength{\tabcolsep}{4pt}
\begin{tabular}{@{}lll@{}}
\toprule
object & here & the causal article \\
\midrule
signal coordinate & $x$ (space) & $t$ (time) \\
transform variable & $\omega$ (Fourier) & $s$ (Laplace) \\
scales in a stage & $r \le s \le t$ & $x \le y \le z$ \\
stage operator and its kernel & $\Phis{s}{t}$, $\mu_{s,t}$ & $\Phi_{x,y}$, $\mu_{x,y}$ \\
exponent of a stage & $g_{s,t}(\omega) = -\log\hat\mu_{s,t}(\omega)$ (Lemma~\ref{lem:nonvanishing}) & $g_{x,y}(s)$ \\
kernel from the origin, accumulated exponent & $\mu_{0,t}$, $G(t,\omega) = g_{0,t}(\omega)$ & $\mu_{0,x}$, $G(x,s)$ \\
smoothed transmittance & $\Theta(t)$ (Corollary~\ref{cor:smoothed-transmittance}) & --- \\
canonical gauge & $\chi(t)$ (Proposition~\ref{prop:canonical-gauge}) & $\chi(x)$ \\
the family's exponent and its pair & $F(\omega)$; Gaussian coefficient $a$, profile $k$ & $F(s)$; drift $b_0$, profile $k$ \\
\Levy{} measure, tail measure of the profile & $\nu$ (\ref{eq:levy-khintchine}), $\varpi = -dk$ (Lemma~\ref{lem:selfdecomposable-exponents}) & $\nu$, $\varpi$ \\
symbol of the exponent & $B(\omega) = \omega F'(\omega)$ (Lemma~\ref{lem:selfdecomposable-exponents}) & $B(s) = sF'(s)$ \\
dilation ratio, action, orbit coordinate & $\lambda$, $S_\lambda$, $c(\lambda)$ & $\sigma$, $S_\sigma$, $c(\sigma)$ \\
random displacement, random delay & $X_{s,t}$ & $T_{x,y}$ \\
function space & $\Lone = L^1(\RR,dx)$ & $X = L^1(\RR,dt)$ \\
translation, reflection & $\Tr{a}$, $R$ & $\tau_a$, --- \\
exponents by definition, by representation & $\NDs$ (Def.~\ref{def:symmetric-negdef}), $\LEs$ (Def.~\ref{def:levy-exponents}) & $\mathrm{BF}_0$, $\mathrm{LE}$ \\
\bottomrule
\end{tabular}
\end{table}

\emph{Spaces.} We write $\Lone = L^1(\RR, dx)$, real, with positive cone $L^1_+$ and
$\int f := \int_\RR f(x)\,dx$. Its dual is $L^\infty(\RR)$. $M_+(\RR)$ denotes the finite
positive Borel measures on $\RR$, of total mass $\|\mu\| = \mu(\RR)$.

\emph{Actions.} Translations are $(\Tr{a} f)(x) = f(x-a)$, $a \in \RR$. The reflection is
$(Rf)(x) = f(-x)$, extended to measures as the pushforward under $x \mapsto -x$, and a
measure is \emph{symmetric} when $R\mu = \mu$. Dilations act on functions by
$(\Dil{\lambda} f)(x) = \lambda^{-1} f(x/\lambda)$, $\lambda > 0$, which preserves mass, and
on measures by the pushforward under $x \mapsto \lambda x$.

\emph{Convolution.} $(\mu * f)(x) = \int f(x-y)\,\mu(dy)$ for $\mu \in M_+(\RR)$ and
$f \in \Lone$. Then $\|\mu * f\|_1 \le \|\mu\|\,\|f\|_1$, and both actions are compatible
with convolution:
\[
  \Dil{\lambda}(\mu * f) = (\Dil{\lambda}\mu) * (\Dil{\lambda}f), \qquad
  R(\mu * f) = (R\mu) * (Rf).
\]

\emph{Transforms.} For $\mu \in M_+(\RR)$ the Fourier transform is
\begin{equation}\label{eq:transform}
  \hat\mu(\omega) = \int_\RR e^{-i\omega x}\,\mu(dx), \qquad \omega \in \RR,
\end{equation}
continuous, with $|\hat\mu| \le \|\mu\|$ and $\hat\mu(0) = \|\mu\|$; for symmetric $\mu$ it is
real and even, $\hat\mu(\omega) = \int \cos(\omega x)\,\mu(dx)$. Under dilation
$\widehat{\Dil{\lambda}\mu}(\omega) = \hat\mu(\lambda\omega)$, and
$\widehat{\mu * \nu} = \hat\mu\,\hat\nu$. Unlike the Laplace transform of a measure carried by
a half-line, $\hat\mu$ need not be positive. That the axioms exclude a zero is
Lemma~\ref{lem:nonvanishing}.

Two classes of functions on $\RR$ carry the analysis: the positive definite functions, which
are the transforms of the kernels, and the symmetric continuous negative definite functions,
which are their exponents.

% shared with blueprint def:positive-definite
\begin{definition}[positive definite]\label{def:positive-definite}
A function $\varphi : \RR \to \CC$ is \emph{positive definite} if
\[
  \sum_{j,k} c_j\,\overline{c_k}\,\varphi(\omega_j - \omega_k) \ge 0
  \quad\text{(in particular real)}
\]
for every finite family $(\omega_j) \subset \RR$ and every finite family $(c_j) \subset \CC$.
\end{definition}

The words in parentheses belong to the definition. The double sum is required to be a
nonnegative real number, and for a complex-valued $\varphi$ that is a genuine condition. It
forces $\varphi$ to be Hermitian, $\varphi(-\omega) = \overline{\varphi(\omega)}$, and it is
the reading under which Proposition~\ref{prop:fourier-toolbox}(1) holds.

% shared with blueprint def:symmetric-negdef
\begin{definition}[symmetric continuous negative definite function]\label{def:symmetric-negdef}
A function $\psi : \RR \to [0,\infty)$ is a \emph{symmetric continuous negative definite
function} if $\psi$ is continuous and even, $\psi(0) = 0$, and $e^{-\tau\psi}$ is positive
definite for every $\tau > 0$. We write $\psi \in \NDs$.
\end{definition}

Definition~\ref{def:symmetric-negdef} defines $\NDs$ by the exponential form, that
$e^{-\tau\psi}$ is positive definite for every $\tau > 0$, and not by the quadratic-form
inequality in $\psi$ that the sources call the kernel form (clause~(2) of the next
proposition). The exponential form is the one the cascade asks for, and a proof that a
given $\psi$ lies in $\NDs$ then has to exhibit the transforms $e^{-\tau\psi}$ and nothing
else. Schoenberg's theorem, the equivalence of the two forms, is needed only to import
results that the sources state in the kernel form. That is why clause~(2) is a cited
interface rather than a definition.

The toolbox that connects the two classes to transforms is classical and is cited rather
than reproven. The sources, clause by clause, follow the statement.

% shared with blueprint prop:fourier-toolbox
\begin{proposition}[the Fourier toolbox]\label{prop:fourier-toolbox}
\begin{enumerate}
\item \emph{(Bochner.)} $\varphi$ is the Fourier transform of a probability measure on $\RR$
if and only if $\varphi$ is continuous, positive definite and $\varphi(0) = 1$; it is the
transform of a \emph{symmetric} probability measure if and only if moreover $\varphi$ is real,
equivalently even.
\item \emph{(Schoenberg.)} For continuous even $\psi \ge 0$ with $\psi(0) = 0$: $\psi \in \NDs$
if and only if $\psi$ is negative definite in the kernel sense,
\[
  \sum_{j,k}\bigl(\psi(\omega_j) + \psi(\omega_k) - \psi(\omega_j - \omega_k)\bigr)\,
  c_j\,\overline{c_k} \ \ge\ 0
\]
for every finite family; equivalently, by clause (1), if and only if $e^{-\tau\psi}$ is the
transform of a symmetric probability measure for every $\tau > 0$.
\item \emph{(\Levy--Khintchine, symmetric form.)} Every $\psi \in \NDs$ has the representation
\begin{equation}\label{eq:levy-khintchine}
  \psi(\omega) \;=\; a\,\omega^2 + \int_{(0,\infty)} \bigl(1 - \cos\omega x\bigr)\,\nu(dx),
  \qquad a \ge 0, \quad \int_{(0,\infty)} (1 \wedge x^2)\,\nu(dx) < \infty,
\end{equation}
the pair $(a,\nu)$ being unique; and conversely every such pair defines a member of $\NDs$.
\item \emph{(Closure.)} $\NDs$ is a convex cone, and in the kernel form of clause (2) the
class is closed under pointwise limits with no proviso. If $\psi_n \in \NDs$ and
$\psi_n \to \psi$ pointwise with $\psi$ continuous at $0$, then $\psi \in \NDs$.
\end{enumerate}
\end{proposition}

\emph{Sources.} Clause~(1) is \sscite{sato1999levy}, Prop.~2.5, pp.~8--9, whose parts are
numbered (i)--(xii): part (i) is Bochner's theorem, and parts (ii) and (v) add that the
measure is unique and that it is symmetric exactly when the transform is real. Feller
corroborates it (\sscite{feller2009introduction}, Vol.~2, \S~XIX.2, Theorem (Bochner),
p.~622). Clause~(2) is \sscite{schilling2012bernstein}, Def.~4.3 and Prop.~4.4, p.~36; also
\sscite{berg1975potential}, Thm.~7.8, p.~41, and \sscite{jacob2001pseudo}, Def.~3.6.5,
p.~122, Def.~3.6.6, p.~123, and Thm.~3.6.11, p.~125. Clause~(3) is \sscite{sato1999levy},
Thm.~8.1, pp.~37--38, whose parts (i), (ii) and (iii) are existence, uniqueness and the
converse, with (8.2) there for the condition on the \Levy{} measure. Clause~(4) is
\sscite{jacob2001pseudo}, Lemma~3.6.7A, p.~123, for the kernel form, and
\sscite{sato1999levy}, Prop.~2.5, part (viii), p.~9, for the proviso. Sato's characteristic
function is $\int e^{i\langle z, x\rangle}\,\mu(dx)$, the opposite sign to
\ref{eq:transform}, which is immaterial for symmetric laws.

% paper-only remark (2026-09-13, at the author's request), condensed from the hub's
% wiki/positive-definite-functions.md sections 1-2; the Feller anchors are those the librarian
% confirmed (Lemma 1 p. 621 uses an even window, the Gaussian factor is his (2.4); the Toeplitz
% form is (2.6)-(2.7) p. 622; Lemma 3 p. 622).
\begin{remark}[what the two classes are, and why the transform sees them]
\label{rem:pd-nd-intuition}
The prototype of a positive definite function is a single \emph{character},
$\omega \mapsto e^{-i\omega x}$ for a fixed $x$: the double sum of
Definition~\ref{def:positive-definite} collapses to $|\sum_j c_j e^{-i\omega_j x}|^2 \ge 0$.
The condition is linear in $\varphi$ and survives pointwise limits, so every positive
mixture of characters, that is every Fourier transform of a positive measure, is positive
definite for free. Bochner's theorem is the converse: a continuous positive definite
function is such a mixture, and the mixing measure is unique. In the language of convex
cones, the characters are the extreme rays of the cone of positive definite functions, and
every continuous member is the barycentre of the rays, with the measure as the weights. In
the symmetric case the rays are $\cos\omega x$, $x \ge 0$.

The negative definite functions are the same picture one level up. Ask which $\psi$ make
$e^{-\tau\psi}$ a transform for every $\tau > 0$, which is what the cascade of
\S\ref{sec:cascade} will ask. The answer is again a cone described twice: by the checkable
kernel-form inequality of clause~(2), and by its extreme rays, which is the
\Levy--Khintchine representation \ref{eq:levy-khintchine}. The rays are $\omega^2$ and
$1 - \cos\omega x$ for $x > 0$, and there is no constant ray because $\psi(0) = 0$. Each ray
has a probabilistic identity, given in Remark~\ref{rem:displacement-dictionary} below.

The mechanism behind Bochner's theorem is the one this article reuses. In Feller's proof the
transform is inverted against an even window and the general case is reached through the
convergence factor $e^{-\varepsilon\omega^2}$ (\sscite{feller2009introduction}, Vol.~2,
\S~XIX.2, Lemma~1, p.~621, and (2.4)); positive definiteness is what makes the windowed
inverse transforms nonnegative, and the finite-family inequality is the discrete form of
that, the Toeplitz form (2.6)--(2.7) on p.~622. The same device, a transform read against a
Gaussian window, is the test function of \S\ref{sec:representation} and the smoothed
transmittance of \S\ref{sec:cascade}. Two facts used throughout are also Feller's
(Lemma~3, p.~622): $|\varphi| \le \varphi(0)$ and $\varphi(-\omega) = \overline{\varphi(\omega)}$.

The contrast with the half-line is one of kind. Complete monotonicity, the transform-free
description of ``comes from positive mass'' for the Laplace transform, is a pointwise test:
countably many derivative signs at each point. Positive definiteness is a global
quadratic-form condition, and there is no local test for it. That is the first sign that
the half-line theory is the more rigid of the two, and the remark at the end of this section
lists what the difference costs.
\end{remark}

\emph{The folding convention.} The sources state clause~(3) on $\RR$, with a centering term
and a \Levy{} measure on $\RR \setminus \{0\}$. The reduction to \ref{eq:levy-khintchine} is
ours and elementary. A symmetric \Levy{} measure $\nu_2$ on $\RR \setminus \{0\}$ enters
\ref{eq:levy-khintchine} as its image $\nu$ under $x \mapsto |x|$. So a two-sided \Levy{}
density written $h(|x|)/|x|$, the form the probabilistic sources use, has folded profile
$k = 2h$ in \ref{eq:sd-profile} below: the factor is the two directions of a displacement
of a given size, and it is the one place where a constant is easily lost in carrying a
statement from a source into this article. The centering term of the general representation
vanishes by symmetry, read on the compensated integrand: the odd part of
$e^{-i\omega x} - 1 + i\omega x\,\mathbf 1_{|x| < 1}$ is $-i(\sin\omega x - \omega x\,\mathbf
1_{|x|<1})$, which is $O(x^3)$ at the origin and bounded, hence $\nu_2$-integrable, and odd,
so its integral against a symmetric $\nu_2$ is $0$. The uncompensated
$\int \sin(\omega x)\,\nu_2(dx)$ need not converge.

\emph{Which vocabulary does the work.} The later sections name the class by its
representation.

% paper-only definition (2026-09-13, at the author's request): the name LE_s is the
% blueprint's and the Lean's, used by chapters 5-7 and glossed there only in prose; a node of
% its own is the recommended blueprint mirror.
\begin{definition}[symmetric \Levy{} exponents]\label{def:levy-exponents}
$\LEs$ is the set of functions $\psi : \RR \to [0,\infty)$ of the form
\ref{eq:levy-khintchine}, for some $a \ge 0$ and some measure $\nu$ on $(0,\infty)$ with
$\int (1 \wedge x^2)\,\nu(dx) < \infty$. By Proposition~\ref{prop:fourier-toolbox}(3),
$\LEs = \NDs$.
\end{definition}

The representation is the object the proofs manipulate, and $\LEs$ is the name the later
sections use. Definition~\ref{def:symmetric-negdef} itself enters only where the class
$\NDs$ is the subject, and in one leg of Lemma~\ref{lem:selfdecomposable-exponents}. The
machine-checked development defines $\LEs$ alone. What that development takes from
Proposition~\ref{prop:fourier-toolbox} is less than the proposition states, and the
difference is part of the trust base stated in \S\ref{sec:machine-checked}. It takes the
forward implication of clause~(1) (a real-valued continuous positive definite function with
value $1$ at the origin is the cosine transform of a symmetric probability measure) and the
converse and
uniqueness clauses of~(3). Clause~(2), the existence half of~(3) and clause~(4) are stated
because they are the toolbox a reader meets, and they are used by no proof: the arguments
below extract the representation of a limit directly (Theorem~\ref{thm:increments-levy}),
so the closure clause is never needed.

% blueprint rem:displacement-dictionary, copied and split in two for the paper (readability
% pass, 2026-09-12): the reading of the representation here, right after the toolbox that
% states it; the log-displacement reading (rem:displacement-profile-log) after the profile form
% it reads. The blueprint's "rate per unit of scale" is not used: the catalogue of one law is an
% intensity, and how the catalogue depends on scale is section 7's result, not a rate.
\begin{remark}[the displacement dictionary]\label{rem:displacement-dictionary}
A symmetric kernel is the law of a \emph{random displacement}: $\Phis{s}{t}f$ reports the
signal at a randomly displaced position $x - X_{s,t}$ with $X_{s,t} \sim \mu_{s,t}$. Its
transform $\hat\mu_{s,t}(\omega) = \Ex\cos(\omega X_{s,t})$ is the contrast that a
sinusoidal probe of wavenumber $\omega$ retains after the random displacement, the stage's
modulation transfer function (Remark~\ref{rem:modulation-transfer}). The exponent
$\psi = -\log\hat\mu$ adds along the cascade. In \ref{eq:levy-khintchine} the Gaussian
coefficient $a$ is the \emph{diffusive part}, a Gaussian displacement of variance $2a$, the
one term the time-causal axis deletes and
the only one Gaussian scale space uses. The measure $\nu$ is a \emph{catalogue of jump
displacements}. Displacements of size about $x$ arrive as a Poisson stream, that is at
independent random moments with a constant rate, and $\nu(dx)$ is that rate; the streams of
different sizes are independent, and each event of size $x$ costs the probe the phase
deficit $1 - \cos\omega x$. So a symmetric infinitely divisible displacement is a Gaussian
part plus an independent sum of jumps of all sizes, and the representation says that every
such displacement decomposes in this way.
\end{remark}

Two further classical facts are used repeatedly and are recorded as statements so that
each later use can cite one. Both are parts of Sato's Prop.~2.5 (\sscite{sato1999levy}):
uniqueness is its part (ii), p.~8, and the continuity theorem its parts (vii)--(viii),
p.~9. Neither is taken on trust: each is proved, in the form the article uses, from a
result in Mathlib, and the proof is given.

% shared with blueprint prop:fourier-uniqueness
\begin{proposition}[uniqueness of the Fourier--Stieltjes transform]\label{prop:fourier-uniqueness}
A finite Borel measure on $\RR$ is determined by its Fourier transform: if
$\hat\mu(\omega) = \hat\rho(\omega)$ for every $\omega \in \RR$, then $\mu = \rho$.
\end{proposition}

\begin{proof}
The functions $e_\omega(x) = e^{-i\omega x}$, $\omega \in \RR$, span a subalgebra $\mathcal{A}$
of the bounded continuous functions on $\RR$: it contains the constants, since $e_0 = 1$; it is
closed under conjugation, since $\overline{e_\omega} = e_{-\omega}$; and it separates points,
since $e_\omega(x) = e_\omega(y)$ for every $\omega$ forces $x = y$. The hypothesis says that
$\mu$ and $\rho$ give equal integrals to every $e_\omega$, hence, by linearity, to every member
of $\mathcal{A}$. Two finite Borel measures on a complete separable metric space that give equal
integrals to every member of such an algebra are equal: on each compact set the
Stone--Weierstrass theorem makes $\mathcal{A}$ uniformly dense in the continuous functions, and a
finite Borel measure on a Polish space is inner regular, so the equality extends first to the
bounded continuous functions and then, by the portmanteau theorem, to the Borel sets. This is
the argument Mathlib's \texttt{MeasureTheory.Measure.ext\_of\_charFun} runs, in the stated
generality; every step of it is machine-checked, and the article's declaration is its
specialisation to $\RR$.
\end{proof}

The statement is about finite measures on $\RR$ and claims nothing for infinite ones.

% shared with blueprint prop:levy-continuity
\begin{proposition}[\Levy's continuity theorem]\label{prop:levy-continuity}
Let $\mu_n, \mu$ be probability measures on $\RR$. If $\hat\mu_n(\omega) \to \hat\mu(\omega)$
for every $\omega \in \RR$, then $\mu_n \to \mu$ weakly.
\end{proposition}

\begin{proof}
Since $\hat\mu$ is continuous at $0$ with $\hat\mu(0) = 1$, the sequence is tight. Indeed, for a
probability measure $\rho$ on $\RR$ and $r > 0$, Fubini's theorem gives
\[
  \frac1r\int_{-r}^{r}\bigl(1 - \hat\rho(\omega)\bigr)\,d\omega
  = 2\int_\RR\Bigl(1 - \frac{\sin rx}{rx}\Bigr)\,\rho(dx)
  \ \ge\ \rho\bigl(\{|x| \ge 2/r\}\bigr),
\]
because the integrand is nonnegative and is at least $\tfrac12$ where $|rx| \ge 2$. Applying
this to $\mu_n$ and letting $n \to \infty$ --- the integrands are bounded by $2$, so bounded
convergence applies --- bounds the upper limit of $\mu_n(\{|x| \ge 2/r\})$ by
$\frac1r\int_{-r}^r(1 - \hat\mu)$, which tends to $0$ as $r \downarrow 0$ by continuity of
$\hat\mu$ at $0$. So $\{\mu_n\}$ is tight. By Prokhorov's theorem every subsequence has a
further subsequence converging weakly to some probability measure $\rho$; along it
$\hat\rho = \lim \hat\mu_n = \hat\mu$, since $x \mapsto e^{-i\omega x}$ is bounded and
continuous, so $\rho = \mu$ by Proposition~\ref{prop:fourier-uniqueness}. Every subsequence
therefore has a further subsequence converging weakly to $\mu$, which is convergence of the
whole sequence. The statement is machine-checked: Mathlib's
\texttt{MeasureTheory.ProbabilityMeasure.tendsto\_of\_tendsto\_charFun} proves it with the same
tightness estimate (\texttt{isTightMeasureSet\_of\_tendsto\_charFun}) and, in place of the
subsequence extraction, a direct density argument in the algebra of characters.
\end{proof}

Weak convergence is read as convergence of $\int g\,d\mu_n$ for every bounded continuous
$g : \RR \to \RR$, which is the form the $\varepsilon/3$ argument in the proof of
Theorem~\ref{thm:main-characterization} uses. What is not part of the statement is the
passage from weak convergence to convergence of the convolution operators in $\Lone$. That is
the $\varepsilon/3$ argument itself, and it is proved where it is used.

% blueprint rem:levy-continuity-general, copied
\begin{remark}[the general form of the continuity theorem]\label{rem:levy-continuity-general}
The classical statement is more general: if $\hat\mu_n \to \varphi$ pointwise with $\varphi$
continuous at $0$, then $\varphi$ is the transform of a probability measure $\mu$ and
$\mu_n \to \mu$ weakly. The tightness estimate in the proof above uses only continuity of the
limit function at the origin, so it applies unchanged. What the general form needs in addition
is that the limit along a Prokhorov subsequence has transform $\varphi$, which identifies
$\varphi$ as a transform. The continuity proviso is what the line costs
against the half-line, where pointwise limits of Laplace exponents on $(0,\infty)$ suffice
--- and it is why
Proposition~\ref{prop:fourier-toolbox}(4) carries it too. No statement of this article uses
the general form: every use of Proposition~\ref{prop:levy-continuity} has the limit law
already in hand.
\end{remark}

The next lemma is a growth bound on the exponents. The literature states it with an
unspecified constant (\sscite{jacob2001pseudo}, Lemma~3.6.22); here it is proved with an
explicit one, in three lines from a split of the integral, and it is used twice later. The
constant $C$ is finite exactly because $\int (1 \wedge x^2)\,\nu < \infty$: the two
integrals inside it are the two halves of that condition.

% shared with blueprint lem:quadratic-growth
\begin{lemma}[continuous negative definite functions grow at most quadratically]
\label{lem:quadratic-growth}
Let $\psi \in \NDs$ have the pair $(a,\nu)$ of \ref{eq:levy-khintchine}. Then
\[
  \psi(\omega) \ \le\ C\,(1 + \omega^2), \qquad \omega \in \RR,
  \qquad C := a + \tfrac12\int_{(0,1]} x^2\,\nu(dx) + 2\,\nu((1,\infty)) \ < \ \infty .
\]
In particular $\psi(\omega) = O(\omega^2)$ as $|\omega| \to \infty$.
\end{lemma}

\begin{proof}
Split the integral in \ref{eq:levy-khintchine} at $x = 1$. For $x \le 1$ use
$1 - \cos\omega x \le \tfrac12\omega^2x^2$; for $x > 1$ use $1 - \cos\omega x \le 2$. This
gives
$\psi(\omega) \le a\omega^2 + \tfrac12\omega^2\int_{(0,1]}x^2\,\nu(dx) + 2\,\nu((1,\infty))
\le C\,(1+\omega^2)$. Finiteness of the two integrals is
$\int (1 \wedge x^2)\,\nu(dx) < \infty$.
\end{proof}

The exponent of a kernel is $-\log$ of its transform, so it vanishes at exactly the
frequencies where the transform equals $1$, where the kernel lets a probe wave through
untouched. Three later proofs turn on where that can happen: the nonvanishing lemma of
\S\ref{sec:representation}, the exclusion of lattice kernels in \S\ref{sec:covariance}, and
the strict monotonicity of the accumulated exponent in \S\ref{sec:characterization}. The
next statement says what the set of such frequencies can be. On a half-line a Laplace
exponent that vanishes at one interior point vanishes identically. On the line the set is
closed under addition and negation and is closed as a set, so it is a closed subgroup of
$\RR$, and a closed subgroup of $\RR$ is $\{0\}$, a lattice $c\ZZ$ with $c > 0$, or all of
$\RR$. The lattice case does occur, and it means that the kernel is carried by a lattice.

% shared with blueprint lem:lattice-zero
\begin{lemma}[the zero set of a symmetric exponent is a closed subgroup]\label{lem:lattice-zero}
Let $\mu$ be a probability measure on $\RR$ and put
$N(\mu) := \{\omega \in \RR : \hat\mu(\omega) = 1\}$. Then $N(\mu)$ is a closed subgroup of
$\RR$, hence $\{0\}$, a lattice $c\ZZ$ with $c > 0$, or $\RR$; the last case holds if and only
if $\mu = \delta_0$. For $\omega_0 \ne 0$,
\[
  \omega_0 \in N(\mu)
  \quad\Longleftrightarrow\quad
  \cos(\omega_0 x) = 1 \ \text{ for } \mu\text{-almost every } x
  \quad\Longleftrightarrow\quad
  \mu \text{ is carried by the lattice } \tfrac{2\pi}{|\omega_0|}\ZZ .
\]
If $\mu$ is symmetric with $\hat\mu = e^{-\psi}$ for some $\psi \in \NDs$, then
$N(\mu) = \{\psi = 0\}$: a symmetric exponent vanishes at a nonzero frequency exactly when its
law is a lattice law. Lattice laws can be infinitely divisible --- the symmetric compound
Poisson law on $\ZZ$, of exponent $c\,(1 - \cos\omega)$, is one --- so membership of $\NDs$
does not exclude them.
\end{lemma}

\begin{proof}
Since $\Rea\bigl(1 - e^{-i\omega x}\bigr) = 1 - \cos\omega x \ge 0$, the identity
$\hat\mu(\omega) = 1$ forces $\int (1 - \cos\omega x)\,\mu(dx) = 0$, hence
$\cos(\omega x) = 1$ for $\mu$-almost every $x$, hence $e^{-i\omega x} = 1$ almost everywhere
and $\hat\mu(\omega) = 1$; the three conditions are therefore equivalent. So $N(\mu)$ is
closed under sums and negation, and closed as a set because $\hat\mu$ is continuous: a closed
subgroup of $\RR$, hence $\{0\}$, some $c\ZZ$ with $c > 0$, or $\RR$. If $N(\mu) = \RR$ then
$\hat\mu \equiv 1 = \hat\delta_0$, so $\mu = \delta_0$ by
Proposition~\ref{prop:fourier-uniqueness}; the converse is immediate. For $\omega_0 \ne 0$,
$\cos(\omega_0 x) = 1$ says $x \in \tfrac{2\pi}{|\omega_0|}\ZZ$, so the almost-everywhere
clause says exactly that $\mu$ is carried by that lattice. When $\hat\mu = e^{-\psi}$ with
$\psi$ real, $\hat\mu(\omega) = 1$ if and only if $\psi(\omega) = 0$. For the example,
$c\,(1-\cos\omega)$ is of the form \ref{eq:levy-khintchine} with $a = 0$ and
$\nu = c\,\delta_1$, hence lies in $\NDs$ by Proposition~\ref{prop:fourier-toolbox}(3); its
law is carried by $\ZZ$, and it is infinitely divisible because $c\,(1-\cos\omega)/n$ is
again of that form.
\end{proof}

The last sentence of the lemma bears on what follows: the exclusion of lattices
cannot come from the function class, so it must come from elsewhere. It comes from
covariance, Lemma~\ref{lem:no-lattice}, and the gauge argument of \S\ref{sec:covariance} is
arranged so as not to need the exclusion at all.

The class this article characterizes is the class $L$ of \sscite{sato1999levy}, Def.~15.1,
p.~90. Its members are infinitely divisible (Prop.~15.5, p.~93). Read at the displacement,
the definition says that every contraction $cX$ of the displacement $X$ is a component of $X$
itself: what is left after shrinking is again a displacement, independent of the shrunk part.

% shared with blueprint def:self-decomposable
\begin{definition}[self-decomposable; class $L$]\label{def:self-decomposable}
A probability measure $\mu$ on $\RR$ is \emph{self-decomposable} if for every $b > 1$ there is
a probability measure $\rho_b$ on $\RR$ with
$\hat\mu(\omega) = \hat\mu(\omega/b)\,\hat\rho_b(\omega)$ for every $\omega$; equivalently, if
$X \seq cX + R_c$ with $R_c$ independent of $X$, for every $c \in (0,1)$, where $X$ has law
$\mu$.
\end{definition}

The classical description of the class is by its \Levy{} measures: \sscite{sato1999levy},
Cor.~15.11, p.~95, the one-dimensional form of Thm.~15.10 there, which also records that
self-decomposability imposes no restriction on the Gaussian part. The corollary is stated
for a general law, with $k$ increasing on $(-\infty,0)$ and decreasing on $(0,\infty)$. The
symmetric specialisation to a profile on $(0,\infty)$ alone is immediate.

% shared with blueprint prop:sd-exponents
\begin{proposition}[symmetric self-decomposable exponents]\label{prop:sd-exponents}
A symmetric infinitely divisible law with pair $(a,\nu)$ is self-decomposable if and only if
its \Levy{} measure has the form $\nu(dx) = k(x)\,x^{-1}dx$ on $(0,\infty)$ with
$k : (0,\infty) \to [0,\infty)$ nonincreasing, so that its exponent is
\begin{equation}\label{eq:sd-profile}
  \psi(\omega) \;=\; a\,\omega^2
  + \int_0^\infty \bigl(1 - \cos\omega x\bigr)\,\frac{k(x)}{x}\,dx,
  \qquad a \ge 0, \quad k \ge 0 \ \text{nonincreasing}.
\end{equation}
Self-decomposability imposes no restriction on the Gaussian coefficient $a$.
\end{proposition}

Proposition~\ref{prop:sd-exponents} is where a reader meets the class, and it is stated for
orientation: it lies on no proof path. The profile form the development needs is proved in
Lemma~\ref{lem:selfdecomposable-exponents}, from the dilation identity, the uniqueness
clause of Proposition~\ref{prop:fourier-toolbox}(3) and a monotone difference-quotient
argument in the log-displacement coordinate, and every later statement cites that lemma and
not this one.

% blueprint rem:displacement-dictionary, second half (see the comment at the first half)
\begin{remark}[the profile in log-displacement]\label{rem:displacement-profile-log}
In the log-displacement coordinate $\theta = \log x$ the density shape of \ref{eq:sd-profile}
is the profile $\tilde k(\theta) = k(e^\theta)$ per unit log-displacement, dilation acts by
translation in $\theta$, and self-decomposability is the profile being nonincreasing: zooming
out may only add displacement intensity. The Gaussian is the profile-free member; the
symmetric stable exponents $|\omega|^\alpha$ are the exactly scale-free profiles
$\tilde k = c\,e^{-\alpha\theta}$. The extreme rays of the cone are the Gaussian ray
$\omega^2$ and the unit steps, whose exponents are the $\Cin$ functions of
Lemma~\ref{lem:cin-rays} (Appendix~\ref{sec:appendix}).
\end{remark}

The integrability conditions in \ref{eq:sd-profile} are the admissibility test every family
in this article is put through, so they are recorded as a statement. Its two halves are the
symmetric forms of the causal criterion: near the origin the integrand is
$\asymp \omega^2x^2\,k(x)/x$, so $\int_0^1 x\,k$ replaces $\int_0^1 k$, while at infinity
$1 - \cos \le 2$ gives the same condition. The criterion takes $k$ measurable and nonnegative
and nothing else, in particular not monotone, because it is later applied to increment
densities $k(u/t) - k(u/s)$, which are not.

% shared with blueprint lem:profile-integrability
\begin{lemma}[the integrability criterion for a profile]\label{lem:profile-integrability}
Let $k : (0,\infty) \to [0,\infty)$ be measurable and put $\nu(dx) := k(x)\,x^{-1}dx$. Then
\[
  \int_{(0,\infty)} (1 \wedge x^2)\,\nu(dx) < \infty
  \qquad\Longleftrightarrow\qquad
  \int_0^1 x\,k(x)\,dx < \infty \ \ \text{and} \ \
  \int_1^\infty \frac{k(x)}{x}\,dx < \infty .
\]
Consequently, by Proposition~\ref{prop:fourier-toolbox}(3), the right-hand pair of conditions
is necessary and sufficient for \ref{eq:sd-profile} to define a member of $\NDs$.
\end{lemma}

\begin{proof}
Both terms below are nonnegative, so
\[
  \int_{(0,\infty)}(1 \wedge x^2)\,\nu(dx)
  = \int_0^1 x^2\,k(x)\,\frac{dx}{x} + \int_1^\infty k(x)\,\frac{dx}{x}
  = \int_0^1 x\,k(x)\,dx + \int_1^\infty \frac{k(x)}{x}\,dx,
\]
and the left side is finite exactly when both terms are. The second sentence is
Proposition~\ref{prop:fourier-toolbox}(3) applied to the pair $(a,\nu)$.
\end{proof}

\emph{Admissible.} A pair $(a,k)$ with $a \ge 0$, $k : (0,\infty) \to [0,\infty)$ nonincreasing
and the two integrals of the criterion finite is called \emph{admissible}, and so is the
exponent \ref{eq:sd-profile} it defines; profiles equal almost everywhere define the same
exponent. The zero exponent, $a = 0$ and $k = 0$, is admissible as an exponent, and it is
the one exponent that (ND) excludes as a family (Theorem~\ref{thm:main-characterization}).

\emph{Classes referred to later.} The \emph{symmetric stable} laws have exponent
$c\,|\omega|^\alpha$ with $0 < \alpha \le 2$ (\sscite{sato1999levy}, Thm.~14.14, p.~86);
\S\ref{sec:characterization} recovers them as the semigroup case. The \emph{variance-gamma}
laws have exponent $\gamma\log(1 + \omega^2/2)$ up to scale, and for $\gamma > \tfrac12$
the Mat\'ern kernels as densities (\sscite{madan1990variance}; \sscite{fischer2025variance}, \S~2.3, for their
self-decomposability); they are the second member of \S\ref{sec:axioms}. Both classes
appear as the answer to ``which laws have this property'', never as tools.

\begin{remark}[relation to the causal theory]\label{rem:causal-relation-preliminaries}
Every object of this section has a twin in the causal article
\sscite{fagerstrom2026hemigroup}, and the correspondence of Table~\ref{tab:notation} is
exact at the level of the classes. $\NDs$ stands where the Bernstein class $\mathrm{BF}_0$
stood, with Bochner's theorem in the role of Bernstein--Widder and Schoenberg's theorem in
the role of ``$e^{-\tau g}$ is completely monotone for every $\tau$''. The representation
\ref{eq:levy-khintchine} is the \Levy--Khintchine form of a Laplace exponent with the
drift slot $b_0\,s$ become the Gaussian slot $a\,\omega^2$ and the same profile $k$ in
\ref{eq:sd-profile}. The criterion of Lemma~\ref{lem:profile-integrability} has
$\int_0^1 x\,k$ where the causal criterion has $\int_0^1 k$. The vanishing lemma and the
closure under limits do not transfer, and on both the line costs more than the half-line.
The causal vanishing lemma, that a Laplace exponent vanishing at one interior point
vanishes identically, has no Fourier analogue (Lemma~\ref{lem:lattice-zero}), and the
exclusion of lattices, free there, is here the work of covariance. Closure under pointwise
limits, which holds for the Bernstein functions on $(0,\infty)$ (though not for the
normalized subclass on $[0,\infty)$, since $1 - e^{-ns}$ converges to the indicator of
$s > 0$), is unconditional here only in the kernel form. For the functions themselves it needs continuity of the limit at the origin,
Proposition~\ref{prop:fourier-toolbox}(4) and Remark~\ref{rem:levy-continuity-general}.
The classical characterization of self-decomposable \Levy{} measures transfers in the
opposite direction. The causal development cites it and uses it in the analysis direction
of its main theorem; here it is Proposition~\ref{prop:sd-exponents}, cited and used
nowhere, since the profile form is proved. At that point the trust base of this article is
the narrower of the two.
\end{remark}
